\begin{textsample}{POS Dim 2 – 2021\_c4 – Score 20.00 – 2021\_c4/enade\_2021\_tec\_gest\_tec\_inf\_2a1ac18999.txt}  \label{ex:f2_pos_033}
CAPÍTULO 4 PERCEPÇÃO DA PROVA Nas análises feitas neste capítulo, constam as percepções dos concluintes da Área de Tecnologia em Gestão da Tecnologia da Informação sobre a prova aplicada no Enade/2021.
Essas percepções foram mensuradas por meio de nove questões que avaliaram desde o grau de dificuldade da prova até o tempo gasto para concluí -la.
As percepções sobre a prova foram relacionadas ao desempenho dos estudantes, à Grande Região de funcionamento do curso, à Categoria Administrativa e ao tipo de Organização Acadêmica da IES.
O \textbf{questionário} de percepção da prova encontra -se ao final do Anexo VIII, onde está a \textbf{reprodução} do exame.
O desempenho dos estudantes foi classificado em quatro quartos.
Para tanto, esse desempenho foi \textbf{ordenado} de forma ascendente.
O percentil 25, P25, também conhecido como primeiro quartil, é a nota de desempenho que deixa um quarto ( 25 \% ) dos valores observados abaixo e três quartos acima.
Na Figura 1, apresenta -se uma ilustração desse conceito.
O quarto inferior de desempenho é composto pelas notas abaixo do primeiro quartil.
Já o percentil 75, P75, também conhecido como terceiro quartil, é o valor para o qual há três quartos ( 75 \% ) dos dados abaixo e um quarto acima dele.
O quarto superior de desempenho é composto pelas notas iguais ou acima do terceiro quartil.
O percentil 50, P50, também conhecido como mediana, é o valor que divide as notas em dois conjuntos de igual tamanho.
O segundo quarto inclui os valores entre o primeiro quartil ( P25 ) e a mediana.
O terceiro quarto contém os valores entre a mediana ( P50 ) e o terceiro quartil ( P75 ).
Vale ressaltar que percentis, quartis e medianas são pontos que, não obrigatoriamente, pertencem ao conjunto original de dados, ao passo que os quartos são \textbf{subconjuntos} dos dados originais.
A seguir, serão apresentados gráficos com resultados selecionados, relativos às nove questões avaliadas por grupos de estudantes.
Nas barras dos gráficos, apresenta -se o percentual de alunos que assinalou uma das opções ou a soma das porcentagens daqueles que assinalaram duas ( ou três ) delas.
Por exemplo, para as questões 1 e 2, nos gráficos, é apresentada a porcentagem total de participantes que assinalou as opções ( D ) Difícil e ( E ) Muito difícil.
Em cada barra, foram assinalados, também, os extremos do intervalo de confiança de 95 \% como linhas verticais unidas por uma linha horizontal na forma da letra H maiúscula.
O \textbf{estimador} de um \textbf{parâmetro} com um certo nível de confiança ( por exemplo, 95,0 \% ) deve conter o \textbf{parâmetro} no intervalo de confiança em 95 \% das vezes.
Na comparação entre os \textbf{estimadores} dos \textbf{parâmetros} de duas classes de uma dada categoria ( por exemplo : Norte e Nordeste nas Grandes Regiões, de primeiro e último quarto dentro de desempenho, de IES Públicas e Privada, ou de Universidades e Faculdades ), associados aos seus respectivos intervalos de confiança, diz -se que não há diferença estatisticamente significativa entre os \textbf{parâmetros} das duas categorias quando há \textbf{interseção} entre os intervalos de confiança, e que há diferença, se os intervalos de confiança são \textbf{disjuntos} ( para maiores informações \textbf{vide} Glossário ).
Nas Tabelas do Anexo II, são apresentados os valores absolutos e a distribuição percentual18 das alternativas válidas das nove questões, segundo o mesmo \textbf{recorte} de desempenho dos alunos e da Grande Região de funcionamento do curso, de Categoria Administrativa e do tipo de Organização Acadêmica da IES. 4.1 GRAU DE DIFICULDADE DA PROVA 4.1.1 COMPONENTE DE FORMAÇÃO GERAL Ao avaliarem “ Qual o grau de dificuldade desta prova na parte de Formação Geral? ” ( Questão 1 ), 32,9 \% do grupo de inscritos e presentes optaram pelas alternativas Difícil ou Muito difícil.
Entretanto, para mais da metade dos estudantes ( 59,6 \% ), o Componente de Formação Geral da prova foi avaliado com grau de dificuldade Médio ( Gráfico 4.1, Gráfico 4.2 e, no Anexo II, a Tabela II.1 ).
O percentual de estudantes que consideraram a prova Difícil ou Muito difícil foi maior na região Norte, onde a proporção foi de 40,0 \%, enquanto a de menor proporção foi a Centro-Oeste, com 31,5 \%.
No Gráfico 4.1, é possível observar que as diferenças entre as regiões não são estatisticamente significativas.
Nas Grandes Regiões, a proporção de presentes à prova que avaliaram o Componente de Formação Geral como de grau de dificuldade Médio esteve entre 56,6 \%, na região Sul, e 60,9 \%, nas regiões Nordeste e Sudeste. 18 Cumpre lembrar uma das convenções para tabelas \textbf{numéricas} ( ao final do \textbf{Sumário} ) sobre a possibilidade de a soma das partes não resultar em 100 \% por questões de \textbf{arredondamento}.
O percentual de alunos que consideraram o componente de Formação Geral da prova Difícil ou Muito difícil foi decrescente em função dos quartos de desempenho : 41,5 \%, no primeiro quarto, e 22,9 \%, no quarto, o grupo de melhor desempenho na prova.
Nos quartos de desempenho intermediários, a proporção de alunos que consideraram a prova Difícil ou Muito difícil foi de 37,4 \%, no segundo quarto, e de 30,4 \%, no terceiro quarto.
As diferenças entre os quartos de desempenho são estatisticamente significativas, com exceção da diferença entre os dois quartos inferiores.
Para todos os quartos de desempenho, a alternativa modal para essa pergunta foi a Médio, com 50,5 \% e 67,0 \% dos respondentes nos quartos extremos, primeiro e quarto, respectivamente.
O percentual de estudantes que consideraram o componente de Formação Geral da prova Difícil ou Muito difícil foi maior entre os estudantes de Faculdades, alcançando 34,1 \%, enquanto entre os estudantes das Centros Universitários, a Organização Acadêmica com a menor incidência, alcançou 29,8 \%.
No Gráfico 4.3, é possível observar que as diferenças entre as Organizações Acadêmicas não são estatisticamente significativas.
Nas Organizações Acadêmicas, a proporção de presentes à prova que avaliaram o Componente de Formação Geral como de grau de dificuldade Médio esteve entre 57,1 \%, nos CEFET/IF, e 62,7 \%, nos Centros Universitários ( ver também a Tabela II.2, no Anexo II ).
O percentual de estudantes que consideraram o componente de Formação Geral da prova Difícil ou Muito difícil foi maior entre os estudantes de cursos de IES Privadas ( 33,4 \% ), sem diferença estatisticamente significativa da proporção de IES Públicas ( 30,0 \% ).
Já a proporção de presentes à prova que avaliaram o Componente de Formação Geral como de grau de dificuldade Médio foi de 61,5 \%, entre os estudantes de IES Públicas, e de 59,4 \%, entre os de IES Privadas ( ver Gráfico 4.4 ). 4.1.2 COMPONENTE DE CONHECIMENTO ESPECÍFICO Ao responderem à Questão 2 – “ Qual o grau de dificuldade desta prova na parte de Componente Específico? ” – 42,0 \% do grupo de estudantes classificaram -na como Difícil ou Muito difícil.
Além disso, o Componente modal de Conhecimento Específico da prova foi considerado de grau de dificuldade Médio por 52,8 \% dos alunos ( Gráfico 4.5, Gráfico 4.6, e, no Anexo II, a Tabela II.3 ).
A análise das respostas dos estudantes quanto ao grau de dificuldade do Componente de Conhecimento Específico da prova, agregado por Grande Região, mostra diferenças estatisticamente significativas entre proporções de alunos que a avaliaram como Difícil ou Muito difícil.
Tais diferenças são observadas entre a região com a menor proporção, a Nordeste ( 32,5 \% ), e as regiões Sudeste ( 41,7 \% ) e Sul ( 44,7 \% ).
O percentual de alunos que classificaram o grau de dificuldade como Médio, no Componente de Conhecimento Específico, variou de 49,3 \% a 60,0 \%, nas regiões Sul e Norte, respectivamente.
Considerando -se a avaliação da dificuldade das questões do Componente de Conhecimento Específico da prova de acordo com o desempenho dos estudantes, observa -se uma tendência decrescente, a partir do segundo quarto, dos que consideraram a parte específica Difícil ou Muito difícil, variando de 38,0 \% ( último quarto ) a 45,7 \% ( segundo quarto ).
As demais proporções foram de 43,7 \% e 40,7 \%, respectivamente, no primeiro e terceiro quartos.
Observa -se que há diferença estatisticamente significativa dos resultados apenas entre o quarto superior de desempenho e o segundo quarto.
Já os que responderam que o grau de dificuldade das questões do Componente de Conhecimento Específico da prova foi Médio atingiram \textbf{percentagens} que variaram de 49,0 \%, no segundo quarto, a 57,2 \%, no último.
Considerando -se a avaliação da dificuldade das questões do Componente de Conhecimento Específico da prova, o percentual de \textbf{estudantes} que as consideraram Difícil ou Muito difícil foi maior entre os estudantes de Universidades, atingindo 43,1 \% ; já entre os de CEFET/IF, a Organização Acadêmica com a menor incidência, essa proporção foi de 32,7 \%.
No Gráfico 4.7, é possível observar que a diferença entre as Organizações Acadêmicas não é estatisticamente significativa.
Nas Organizações Acadêmicas, a proporção de presentes à prova que consideraram o Componente de Conhecimento Específico de grau de dificuldade Médio esteve entre 51,5 \%, nas Universidades, e 60,2 \%, nos CEFET/IF ( ver também Tabela II.4, no Anexo II ).
O percentual de estudantes que consideraram o Componente Específico da prova Difícil ou Muito difícil foi maior entre os estudantes de cursos de IES Privadas ( 42,7 \% ), sem diferença estatisticamente significativa da proporção registrada em IES Públicas ( 37,3 \% ).
Já a proporção de presentes à prova que avaliaram esse componente como de grau de dificuldade Médio foi de 58,6 \%, entre os estudantes de IES Públicas, e de 51,9 \%, entre os de IES Privadas ( ver Gráfico 4.8 e Tabela II.4, no Anexo II ). 4.2 EXTENSÃO DA PROVA EM RELAÇÃO AO TEMPO TOTAL \textbf{Indagados} quanto à extensão da prova, em relação ao tempo total oferecido para a sua resolução ( Questão 3 ), os estudantes apontaram, com maior incidência, a alternativa que considerava a extensão Adequada, em todas as agregações consideradas ( Gráfico 4.9, Gráfico 4.10, e, no Anexo II, a Tabela II.5 ).
O percentual de alunos que responderam ser a extensão da prova Adequada foi de 63,7 \%.
Já 30,6 \% dos inscritos presentes consideraram que a prova foi Longa ou Muito longa, e 5,7 \% a avaliaram como Curta ou Muito curta.
Entre as Grandes Regiões, os que avaliaram a prova como Longa ou Muito longa, em relação ao tempo total destinado à sua resolução, atingiram uma porcentagem que variou de 28,1 \%, na região Nordeste, até 33,8 \%, na região Centro-Oeste.
Não há diferença estatisticamente significativa entre as regiões.
Considerando -se o desempenho dos alunos, nota -se ainda que 63,3 \% consideraram a extensão da prova Adequada, no quarto de desempenho inferior, e 64,5 \%, no de melhor desempenho ( quarto superior ).
Nos quartos intermediários, essa proporção foi de 61,2 \%, no segundo quarto, e de 65,7 \%, no terceiro.
No Gráfico 4.10, observa -se que as proporções de estudantes que consideraram a prova Longa ou Muito longa não apresentam uma tendência específica em função dos quartos de desempenho, sendo a maior proporção dessa resposta no segundo quarto de desempenho ( 31,6 \% ), e a menor, no terceiro quarto ( 30,1 \% ).
Pode -se constatar que não há diferença estatisticamente significativa entre os quartos de desempenho.
Considerando -se a extensão da prova, o percentual de estudantes que avaliaram a prova como Longa ou Muito longa, em relação ao tempo total destinado à sua resolução, foi maior entre os estudantes de CEFET/IF, alcançando 31,6 \%, enquanto entre os estudantes de Centros Universitários, a Organização Acadêmica com a menor incidência, esse índice atingiu 29,5 \%.
No Gráfico 4.11, é possível observar que não há diferença estatisticamente significativa entre os tipos de Organização Acadêmica.
A proporção de presentes à prova que consideraram a sua extensão Adequada, em relação ao tempo total destinado à sua resolução, esteve entre 61,2 \%, nos CEFET/IF, e 65,1 \%, nos Centros Universitários ( ver também Tabela II.6, no Anexo II ).
O percentual de estudantes que consideraram a prova Longa ou Muito longa, em relação ao tempo total destinado à sua resolução, foi maior entre os estudantes de cursos de IES Privadas ( 30,4 \% ), não havendo diferença estatisticamente significativa da proporção em relação aos de IES Públicas ( 31,5 \% ).
Já a proporção de presentes à prova que consideraram a extensão Adequada foi de 65,0 \%, entre os estudantes de IES Públicas, e de 63,5 \%, entre os de IES Privadas ( ver Gráfico 4.12 e Tabela II.6, no Anexo II ). 4.3 COMPREENSÃO DOS ENUNCIADOS DAS QUESTÕES 4.3.1 COMPONENTE DE FORMAÇÃO GERAL Em relação aos enunciados das questões do Componente de Formação Geral ( Questão 4 ), as opiniões foram positivas, já que 75,1 \% dos alunos avaliados consideraram todos ou a maioria dos enunciados das questões “ claros e objetivos ” ( Gráfico 4.13, Gráfico 4.14 e, no Anexo II, a Tabela II.7 ).
Na análise regional, a \textbf{percentagem} de estudantes que avaliaram que todos ou a maioria dos enunciados das questões do Componente de Formação Geral estavam “ claros e objetivos ” variou de 70,0 \%, na região Norte, a 77,9 \%, na região Sul.
A análise das percepções dos estudantes sobre a clareza e a objetividade dos enunciados permite afirmar que todos, ou a maioria dos enunciados de questões relativas ao Componente de Formação Geral, foram considerados “ claros e objetivos ” para a maior parte dos respondentes ( maior ou igual a 70,0 \%, em todas as regiões, e maior ou igual a 66,2 \%, em todos os quartos de desempenho ).
Segundo o desempenho, observa -se que a proporção dos que emitiram essa opinião \textbf{cresce}, conforme o desempenho aumenta, com diferenças estatisticamente significativas entre os quartos inferiores de desempenho e os quartos superiores.
No quarto superior, a clareza e a objetividade de todos ou da maioria dos enunciados das questões foram percebidas por 82,5 \%.
A proporção dos presentes que consideraram que todos ou a maioria dos enunciados das questões do Componente de Formação Geral estavam “ claros e objetivos ” foi maior entre os estudantes de Centros Universitários, chegando a 77,7 \%, enquanto entre os de Universidades, a Organização Acadêmica com a menor incidência, esse índice atingiu 73,9 \%.
No Gráfico 4.15, é possível observar que as diferenças entre os tipos de Organização Acadêmica não são estatisticamente significativas ( ver também Tabela II.8, no Anexo II ).
O percentual de estudantes que consideraram que todos ou a maioria dos enunciados das questões do Componente de Formação Geral estavam “ claros e objetivos ” foi maior entre os estudantes de cursos de IES Públicas ( 76,1 \% ), mas sem diferença estatisticamente significativa da proporção registrada em IES Privadas ( 75,0 \% ).
Um maior detalhamento pode ser encontrado no Gráfico 4.16 e na Tabela II.8, no Anexo II. 4.3.2 COMPONENTE DE CONHECIMENTO ESPECÍFICO Em relação aos enunciados das questões do Componente de Conhecimento Específico da prova, para 76,1 \% dos estudantes avaliados da Área de Tecnologia em Gestão da Tecnologia da Informação, a clareza e a objetividade ( Questão 5 ) estavam presentes em todas ou na maioria das questões ( Gráfico 4.17 a Gráfico 4.20, e no Anexo II, as Tabelas II.9 e II.10 ).
A maioria dos estudantes de todas as Grandes Regiões brasileiras considerou “ claros e objetivos ” todos ou a maioria dos enunciados das questões do Componente de Conhecimento Específico da prova, percentual sempre maior ou igual a 50,0 \%.
A diferença entre as regiões Sudeste e Sul é a única estatisticamente significativa.
A proporção de estudantes que consideraram os enunciados das questões “ claros e objetivos ” apresenta uma tendência crescente em relação ao aumento de desempenho : mais elevada no quarto superior ( 84,6 \% ), se comparada ao quarto inferior de desempenho ( 67,5 \% ).
As diferenças entre os quartos inferiores de desempenho e os quartos superiores são estatisticamente significativas.
A proporção dos presentes que consideraram que todos ou a maioria dos enunciados das questões do Componente de Conhecimento Específico estavam “ claros e objetivos ” foi maior entre os estudantes de Centros Universitários, atingindo 78,8 \%.
Entre os estudantes de Universidades, a Organização Acadêmica com a menor incidência, essa proporção atingiu 74,6 \%.
No Gráfico 4.19, é possível observar que as diferenças entre os tipos de Organização Acadêmica não são estatisticamente significativas ( ver também Tabela II.10, no Anexo II ).
O percentual de estudantes que consideraram que todos ou a maioria dos enunciados das questões do Componente de Conhecimento Específico estavam “ claros e objetivos ” foi maior entre os estudantes de cursos de IES Públicas ( 78,5 \% ), sem diferença estatisticamente significativa da proporção registrada em IES Privadas ( 75,7 \% ).
Ver Gráfico 4.20 e Tabela II.10, no Anexo II, para maior detalhamento. 4.4 SUFICIÊNCIA DAS INFORMAÇÕES/INSTRUÇÕES FORNECIDAS Ao se avaliarem as informações/instruções fornecidas para a resolução das questões ( Questão 6 ), constatou -se que 82,9 \% dos respondentes da Área de Tecnologia em Gestão da Tecnologia da Informação de todo o Brasil afirmaram que eram até excessivas ou suficientes em todas ou na maioria das questões ( Gráfico 4.21, Gráfico 4.22 e, no Anexo II, a Tabela II.11 ).
Quanto à distribuição de respondentes pelas Grandes Regiões, observa -se que a proporção de estudantes que consideraram as informações/instruções fornecidas até excessivas ou suficientes em todas ou na maioria das questões foi sempre superior ou igual a 79,2 \%, chegando a 90,0 \%, na região Norte.
Não existem diferenças estatisticamente significativas entre as regiões.
Levando -se em conta o desempenho dos participantes, observa -se uma tendência crescente das proporções de participantes que avaliaram as informações/instruções como até excessivas ou suficientes em todas ou na maioria das questões, havendo diferenças estatisticamente significativas entre os quartos de desempenho, com exceção da diferença entre os quartos inferiores, conforme mostrado no Gráfico 4.22.
O percentual foi mais elevado no quarto superior ( 89,7 \% ), percentual superior à média nacional ( 82,9 \% ).
Já no quarto inferior, a suficiência das informações/instruções declarada como até excessiva, em todas ou na maioria das questões foi percebida por 77,4 \% dos respondentes.
Considerando -se as informações/instruções fornecidas para a resolução das questões ( Questão 6 ), a proporção que considerou que eram até excessivas ou suficientes em todas ou na maioria das questões foi maior entre os estudantes de Centros Universitários, \textbf{totalizando} 87,6 \%.
Os CEFET/IF foram a Organização Acadêmica em que houve a menor proporção : 77,6 \%.
No Gráfico 4.23, é possível observar que as diferenças entre Centros Universitários em relação às Universidades ( 81,8 \% ) e Faculdades ( 80,7 \% ), são estatisticamente significativas ( ver também Tabela II.12, no Anexo II ).
O percentual de estudantes que consideraram que as informações/instruções fornecidas para a resolução das questões eram até excessivas ou suficientes em todas ou na maioria das questões foi menor entre os estudantes de cursos de IES Públicas ( 81,6 \% ), não havendo diferença estatisticamente significativa da proporção registrada em IES Privadas ( 83,1 \% ) ( ver Gráfico 4.24 e Tabela II.12, no Anexo II ). 4.5 DIFICULDADE ENCONTRADA AO RESPONDER À PROVA \textbf{Indagados} sobre as dificuldades com as quais se depararam ao responder à prova ( Questão 7 ), 14,0 \% dos estudantes apontaram o Desconhecimento do conteúdo.
Para 43,0 \%, a Forma diferente de abordagem do conteúdo foi indicada como dificuldade.
Já a Falta de motivação para fazer a prova foi a dificuldade apontada por 17,4 \% dos respondentes.
Considerando -se todo o Brasil, 21,7 \% dos respondentes afirmaram que não tiveram qualquer tipo de dificuldade para responder à prova ( Tabelas II.13 e II.14, no Anexo II ).
Nos Gráficos 4.25 a 4.28, são apresentados os percentuais de estudantes que apontaram o Desconhecimento do conteúdo como “ dificuldade ao responder à prova ”.
Na análise por Grandes Regiões, o percentual de inscritos e presentes que apontaram o Desconhecimento do conteúdo como “ dificuldade ao responder à prova ” não superou 30,0 \%.
Os percentuais variaram de 13,3 \%, na região Sul, a 30,0 \%, na região Norte.
A Forma diferente de abordagem do conteúdo foi a escolha modal dos estudantes, com percentuais que variaram de 30,0 \% ( região Norte ) a 43,4 \% ( região Sudeste ).
O percentual de alunos que citaram a Falta de motivação como dificuldade variou de 10,0 \% ( região Norte ) a 19,4 \% ( região Centro-Oeste ).
Os que declararam não ter qualquer tipo de dificuldade para responder à prova variaram de 21,1 \%, na região Sudeste, a 30,0 \%, na região Norte.
Em relação aos quartos de desempenho, o Desconhecimento do conteúdo foi a opção escolhida por 15,6 \% dos estudantes do quarto inferior e por 13,3 \% do quarto superior, caracterizando uma tendência decrescente.
Não são observadas diferenças estatisticamente significativas entre os quartos de desempenho.
A alternativa modal entre os alunos, quando agregados pelos quartos de desempenho, para a dificuldade encontrada, foi causada pela Forma diferente de abordagem do conteúdo : 40,1 \%, no primeiro quarto, e 45,5 \% do segundo quarto, escolheram essa alternativa.
Na análise por Organização Acadêmica, o percentual de inscritos e presentes que apontou o Desconhecimento do conteúdo como “ dificuldade ao responder à prova ” não superou 25,3 \%.
Os percentuais variaram de 12,8 \%, nas Universidades, a 25,3 \%, nos CEFET/IF.
A Forma diferente de abordagem do conteúdo foi a escolha modal dos estudantes, com percentuais que variaram de 40,4 \% ( CEFET/IF ) a 45,0 \% ( Faculdades ).
O percentual de alunos que citaram a Falta de motivação como dificuldade variou de 15,2 \% ( CEFET/IF ) a 17,9 \% ( Universidades ).
Os que declararam não ter qualquer tipo de dificuldade para responder à prova variaram de 15,2 \%, nos CEFET/IF, a 23,8 \%, nos Centros Universitários.
No Gráfico 4.27, é possível observar que há diferenças estatisticamente significativas entre os CEFET/IF, em relação aos Centros Universitários ( 14,0 \% ) e às Universidades ( ver também Tabela II.14, no Anexo II ).
Em relação às Categorias Administrativas, o Desconhecimento do conteúdo foi a opção escolhida por 17,4 \% dos estudantes, nas IES Públicas, e por 13,6 \%, nas IES Privadas.
Não são observadas diferenças estatisticamente significativas entre essas duas categorias.
A alternativa modal entre os alunos, quando agregados pelas Categorias Administrativas, para a dificuldade encontrada, foi causada pela Forma diferente de abordagem do conteúdo : 45,7 \%, entre os estudantes de IES Públicas, e 42,6 \%, entre os de IES Privadas ( ver Gráfico 4.28 e Tabela II.14, no Anexo II ). 4.6 CONTEÚDOS DAS QUESTÕES OBJETIVAS DA PROVA Ao se analisarem os conteúdos das questões objetivas da prova ( Questão 8 ), verifica -se que um percentual muito pequeno dos estudantes avaliados, apenas 4,7 \%, afirmou que não estudou ainda a maioria desses conteúdos ( Gráficos 4.29 a Gráfico 4.32, e as Tabelas II.15 e II.16, no Anexo II ).
A maioria ( 62,7 \% ) afirmou ter estudado e aprendido muitos ou todos os conteúdos avaliados.
Na análise por Grande Região, a proporção de respondentes que escolheram a opção não estudou ainda a maioria desses conteúdos foi pequena.
Observa -se que, nas regiões Nordeste ( 6,1 \% ) e Centro-Oeste ( 8,5 \% ), as proporções foram ligeiramente maiores do que a média nacional ( 4,7 \% ).
Nas regiões Norte, Sudeste e Sul, a situação foi \textbf{inversa}, com proporções mais baixas, correspondendo a, respectivamente, 0,0 \%, 4,6 \% e 4,3 \%.
Desconsiderando -se a proporção nula da região Norte, não são observadas diferenças estatisticamente significativas entre as regiões.
Em todas as regiões, a maioria dos presentes afirmou ter estudado e aprendido muitos ou todos os conteúdos, com proporções variando entre 60,0 \%, nas regiões Norte e Centro-Oeste, a 67,5 \%, nas regiões Nordeste e Sul.
Considerando -se separadamente as opiniões de estudantes dos quatro quartos de desempenho, observa -se que, no quarto inferior, 9,1 \% ofereceram como resposta que não estudou ainda a maioria desses conteúdos, sendo 1,3 \% os do quarto superior com a mesma resposta.
A diferença entre os alunos que optaram por esse motivo de dificuldade nos quartos de desempenho é estatisticamente significativa.
Tendo -se em conta o quarto superior, 78,3 \% dos alunos afirmaram ter estudado e aprendido muitos ou todos os conteúdos.
No outro extremo, no primeiro quarto, 48,1 \% optaram pelas mesmas categorias.
As proporções são crescentes com o desempenho.
Na análise por Organização Acadêmica, a proporção de respondentes que escolheram a opção não estudou ainda a maioria desses conteúdos foi pequena.
Observa -se que a proporção nas Universidades, com 5,4 \%, foi maior do que a média nacional ( 4,7 \% ).
Nos Centros Universitários ( 3,8 \% ), nas Faculdades ( 3,9 \% ) e nos CEFET/IF ( 4,0 \% ), a situação foi \textbf{inversa} : a proporção foi mais baixa.
No Gráfico 4.31, é possível observar que as diferenças entre os tipos de Organização Acadêmica não são estatisticamente significativas.
Em todos os tipos de Organização Acadêmica, a maioria dos presentes afirmou ter estudado e aprendido muitos ou todos os conteúdos, com proporções variando entre 59,7 \%, nas Faculdades, a 74,7 \%, nos CEFET/IF ( ver também Tabela II.16, no Anexo II ).
Em relação às Categorias Administrativas, observa -se que, nas IES Públicas, 2,7 \% ofereceram como resposta que não estudou ainda a maioria desses conteúdos, sendo 5,0 \% nas IES Privadas com a mesma resposta.
A diferença entre os alunos que optaram por esse motivo de dificuldade nas duas categorias é estatisticamente significativa.
Considerando -se as IES Públicas, 66,5 \% dos alunos afirmaram ter estudado e aprendido muitos ou todos os conteúdos, sendo de 62,1 \% a proporção dos que escolheram essa resposta nas IES Privadas ( ver Gráfico 4.32 e Tabela II.16, no Anexo II ). 4.7 TEMPO GASTO PARA CONCLUIR A PROVA Ao responderem sobre o tempo de conclusão da prova ( Questão 9 ), mais de dois terços dos estudantes ( 69,7 \% ) afirmaram ter gasto Entre duas e quatro horas ( Gráfico 4.33 a Gráfico 4.36 e, no Anexo II, as Tabelas II.17 e II.18 ).
Considerando -se as cinco Grandes Regiões brasileiras, os percentuais dos que utilizaram Entre duas e quatro horas para finalizar a prova nas regiões Sudeste ( 69,5 \% ) e Sul ( 67,6 \% ) foram inferiores ao percentual nacional.
Nas demais Grandes Regiões, os percentuais de alunos que \textbf{dispensaram} Entre duas e quatro horas para concluir a prova foram os seguintes : 80,0 \% na região Norte, 78,5 \% na região Nordeste e 75,0 \%, na região Centro-Oeste, como mostrado no Gráfico 4.33.
Há diferenças estatisticamente significativas entre a região Nordeste e as regiões Sudeste e Sul.
Uma vez analisadas às alternativas escolhidas pelos estudantes que se situam nos diferentes quartos de desempenho, observa -se uma tendência crescente : uma proporção maior de participantes no quarto superior declarou ter gasto Entre duas e quatro horas para concluir a prova quando comparada com a proporção de participantes nos quartos inferiores.
Há diferenças estatisticamente significativas entre o quarto inferior e os demais quartos de desempenho, bem como entre o segundo e o quarto superior de desempenho.
Na análise por Organização Acadêmica, os percentuais dos que utilizaram Entre duas e quatro horas para finalizar a prova nas Universidades ( 69,6 \% ) e Centros Universitários ( 68,4 \% ) foram inferiores ao percentual nacional ( 69,7 \% ).
Nas Faculdades e CEFET/IF, tal percentual foi, respectivamente, de 70,1 \% e 79,8 \%, e, como mostrado no Gráfico 4.35, há diferença estatisticamente significativa entre as proporções dos CEFET/IF em relação às Universidades e aos Centros Universitários.
Em relação às Categorias Administrativas, observa -se que, nas IES Públicas, 70,2 \% declararam ter gasto Entre duas e quatro horas para concluir a prova, sendo 69,6 \% nas IES Privadas com a mesma resposta.
A diferença entre os alunos que optaram por esse motivo de dificuldade nas duas categorias não é estatisticamente significativa ( ver Gráfico 4.36 e Tabela II.18, no Anexo II ).

% matched lemmas: arredondamento, crescer, disjunto, dispensar, estimador, estudantes, indagar, interseção, inverso, numérico, ordenar, parâmetro, percentagem, questionário, recorte, reprodução, subconjunto, sumário, totalizar, vidar
\end{textsample}
